%%%%%%%%%%%%%%%%%%%% chapter.tex %%%%%%%%%%%%%%%%%%%%%%%%%%%%%%%%%%
% Chapter 12: Domain-Specific NoCs for LLMs, Transformers & Neuromorphic SNNs
% Springer Nature Monograph svmono format
%%%%%%%%%%%%%%%%%%%%%%%%%%%%%%%%%%%%%%%%%%%%%%%%%%%%%%%%%%%%%%%%%%%

\chapter{Domain-Specific NoCs for LLMs, Transformers \& Neuromorphic SNNs}
\label{chap:ai-llm-snn}

\abstract*{In the first two decades of Network-on-Chip research, on-chip communication fabrics were engineered for symmetric multi-core processors running cache-coherent CPU hierarchies and stochastic, uniform traffic. Today, computing is defined by \textbf{Generative Artificial Intelligence (AI)} and \textbf{Bio-Inspired Neuromorphic Architectures}. Large Language Models (LLMs) featuring hundreds of billions of parameters (e.g., Llama, GPT, DeepSeek) and event-driven Spiking Neural Networks (SNNs) impose communication patterns fundamentally alien to traditional meshes: massive synchronous All-Reduce and All-to-All collectives across thousands of tensor cores, multi-terabyte KV-cache memory bandwidth bottlenecks, and high-fanout sparse spike multicast trees. In this chapter, we explore domain-specific interconnect innovations designed to overcome these scaling barriers: In-Network Computing where router crossbars execute gradient accumulation in-flight (PAC-NoC), wafer-scale megacore engines and Ruche wire-maximal networks (the 2,048-core 12-nm HammerBlade SoC), 3D face-to-face copper-pillar hybrid bonding providing $>16\,\mathrm{TB/s}$ of vertical DRAM bandwidth (SiloBreaker), and hypergraph spike-sharing partitioning algorithms (SNNcut and SparrowSNN).}

\section{The Generative AI Communication Crisis}
\label{sec:ai-interconnect-intro}

\begin{svgraybox}
\textbf{Architectural Intuition: The Classical Highway vs. The AI Assembly Line} \\
Why do modern Large Language Models and neuromorphic networks break traditional Networks-on-Chip?
\begin{itemize}
    \item \textbf{Classical CPU Traffic (City Traffic)}: In a traditional multi-core CPU, individual cores send small, unpredictable cache-line requests back and forth—like cars navigating street corners in a busy downtown. Standard 2D meshes handle this stochastic traffic adequately.
    \item \textbf{Generative AI \& LLMs (The Industrial Freight Train)}: Training and serving a 100-Billion Parameter Transformer is not random city traffic. Every fraction of a second, thousands of tensor cores simultaneously halt computation and demand to exchange billions of floating-point numbers in a synchronized \textbf{All-Reduce collective}! If packets must travel across dozens of router hops and round-trip through central DRAM, the entire supercomputer stalls in massive pipeline bubbles~\cite{Ouyang2026PAC}.
    \item \textbf{Neuromorphic SNNs (The Forest Fire Alarm)}: Spiking brain models do not stream continuous dense matrices. Instead, an artificial neuron stays dead silent until its membrane voltage crosses a threshold, then flashes a single binary spike that must be duplicated and delivered to hundreds of distant target neurons in a single clock cycle!
\end{itemize}
\end{svgraybox}

In modern distributed Large Language Model architectures, multi-die parallelism is partitioned across three orthogonal architectural axes:

\begin{enumerate}
    \item \textbf{Tensor Parallelism (TP)}: Individual General Matrix Multiply (GEMM) operations within multi-head attention and feed-forward networks (FFN) are split column-wise and row-wise across $P_{\mathrm{TP}}$ cores. At the end of every Transformer sub-layer, an \textbf{All-Reduce} collective is mandatory to sum partial activations:
    \begin{equation}
    V_{\mathrm{AllReduce}} = 2 \left( \frac{P_{\mathrm{TP}} - 1}{P_{\mathrm{TP}}} \right) M \quad \text{bytes}
    \label{eq:allreduce-volume}
    \end{equation}
    where $M$ is the tensor buffer size.
    
    \item \textbf{Pipeline Parallelism (PP)}: Successive layers of the deep network are partitioned across sequential pipeline stages. When micro-batches enter the pipeline, compute bubbles appear during pipeline fill and drain phases:
    \begin{equation}
    F_{\mathrm{bubble}} = \frac{P_{\mathrm{PP}} - 1}{N_{\mathrm{micro}} + P_{\mathrm{PP}} - 1}
    \end{equation}
    Minimizing bubble time requires ultra-low latency point-to-point activations passing across stage boundaries.
    
    \item \textbf{Mixture-of-Experts (MoE)}: Top-$k$ gating routers dynamically dispatch input tokens to specialized expert compute tiles. This creates bursty, unpredictable \textbf{All-to-All} personalized communication patterns that saturate bisection bandwidth and trigger severe crossbar buffer hotspots.
\end{enumerate}

\section{In-Network Computing \& Hierarchical Reduction Trees}
\label{sec:in-network-computing}

In conventional on-chip networks, performing an All-Reduce requires streaming all unreduced gradient flits from source cores across the mesh to central memory or a master root tile, performing summation, and broadcasting the results back. This incurs $2N$ global transfers and quickly saturates the bisection bandwidth.

In contrast, \textbf{In-Network Computing} integrates arithmetic reduction logic directly into router crossbars~\cite{Ouyang2026PAC}:

\subsection{Mathematical Reduction Formulation}
Let $\mathcal{N}_{\mathrm{in}}$ denote the set of input ports participating in a synchronized collective reduction with matching collective tag ID. In a 4-ary aggregation tree, rather than arbitrating for consecutive egress crossbar passes, an internal arithmetic unit combines incoming operand vectors in-flight:
\begin{equation}
\mathbf{V}_{\mathrm{accum}} = \sum_{p \in \mathcal{N}_{\mathrm{in}}} \mathbf{V}_p = \mathbf{V}_{\mathrm{North}} + \mathbf{V}_{\mathrm{East}} + \mathbf{V}_{\mathrm{South}} + \mathbf{V}_{\mathrm{West}}
\label{eq:in-network-sum}
\end{equation}
By computing partial reductions at each router hop:
\begin{itemize}
    \item The total number of transmitted flits is reduced by up to $75\%$ ($1 - 1/K$ for $K$-ary input ports).
    \item Global network contention drops from $O(N \cdot M)$ to $O(M \log N)$, slashing attention backward pass latency by more than $2.4\times$~\cite{Ouyang2026PAC}.
\end{itemize}

\subsection{Hardware Floating-Point Reduction Unit}
To support modern mixed-precision deep learning, the reduction router integrates IEEE 754-compatible FP16, BF16, and INT8 multi-operand adder trees into the switch allocation stage. Flits carry a 4-bit \texttt{reduction\_op} field specifying sum, max, or min aggregation, while lock-step barrier credit counters ensure flits from all participating branches arrive before reduction proceeds.

\section{Wafer-Scale Engines \& Ruche Wire-Maximal Topologies}
\label{sec:wafer-scale-ruche}

Scaling AI accelerators to tens of thousands of tensor cores requires transcending single-die lithographic reticle boundaries ($\approx 858\,\mathrm{mm}^2$):

\subsection{Wafer-Scale Interconnects (Theseus)}
By patterning entire $300\,\mathrm{mm}$ silicon wafers as continuous computation fabrics, cross-reticle communication replaces high-loss off-package PCIe cables with on-wafer metallization tracks~\cite{Zhu2025}:
\begin{itemize}
    \item \textbf{Defect-Tolerant Bypass Switching}: Silicon manufacturing defects inevitably destroy individual router tiles across a wafer. Hardware bypass multiplexers dynamically stitch functional cores into 2D meshes without software recompilation.
    \item \textbf{Sub-Nanosecond Inter-Reticle Latency}: Transmission across scribe lines achieves single-cycle propagation, enabling wafer-scale All-to-All collective bandwidth exceeding $100\,\mathrm{TB/s}$.
\end{itemize}

\subsection{Ruche Networks in Silicon (HammerBlade)}
In the 12-nm FinFET 2,048-core HammerBlade SoC, researchers deployed \textbf{Ruche Networks}~\cite{Gao2026}. 

\begin{svgraybox}
\textbf{The Wire Hierarchy Physics of Ruche Links} \\
In deep-submicron CMOS, metallization layers have vastly different physical cross-sections:
\begin{itemize}
    \item \textbf{Lower Metals (M1--M3)}: Extremely thin ($w \approx 40\,\mathrm{nm}$), high sheet resistance ($R_{\square} > 0.5\,\Omega/\square$), ideal for local standard-cell connections.
    \item \textbf{Upper Metals (M7--M9)}: Massive copper tracks ($w \approx 400\text{--}800\,\mathrm{nm}$), low sheet resistance ($R_{\square} \approx 0.03\,\Omega/\square$).
\end{itemize}
Ruche networks exploit this physical asymmetry by routing long express links across upper metal layers. These express wires hop over 2 or 4 intermediate routers in a single clock cycle \textbf{without adding ports} to intermediate crossbars!
\end{svgraybox}

By eliminating intermediate router pipeline stages, a Ruche-$X$ network reduces network diameter from $D_{\mathrm{mesh}} = 2\sqrt{N}$ down to:
\begin{equation}
D_{\mathrm{ruche}} = \frac{2\sqrt{N}}{X}
\end{equation}
For $X = 4$ in a 64-core cluster, diameter drops from 16 down to 4 hops, doubling collective throughput under equal silicon area.

\section{3D Face-to-Face Hybrid Bonding \& Memory Silo Breaking}
\label{sec:hybrid-bonding-silos}

Because auto-regressive LLM token generation (the decode phase) evaluates single tokens sequentially, performance is bound by memory bandwidth rather than compute throughput. Every token requires streaming all layer weights ($2\,\mathrm{bytes/param}$) and KV-cache tensors:
\begin{equation}
\mathrm{BW}_{\mathrm{required}} = \mathrm{Model\_Size} \times \mathrm{Tokens/sec}
\end{equation}
For a 70-Billion parameter model serving $100\,\mathrm{tokens/s}$, required memory bandwidth exceeds $14\,\mathrm{TB/s}$—far exceeding the $3.2\,\mathrm{TB/s}$ limit of 8-stack HBM3!

In \textbf{SiloBreaker}~\cite{Wu2026Silo}, logic dies and high-density DRAM dies are assembled using \textbf{face-to-face hybrid bonding}:
\begin{itemize}
    \item \textbf{Direct Copper-to-Copper Bonding}: Bond pitch is scaled to $3\,\mu\mathrm{m}$, achieving interconnect density of $>110{,}000\,\mathrm{pads/mm^2}$ ($100\times$ denser than TSV micro-bumps).
    \item \textbf{Vertical Memory Silos}: Each router tile is directly connected to an overhead DRAM column, delivering over $16\,\mathrm{TB/s}$ of vertical memory bandwidth directly into tensor register files.
    \item \textbf{Architecture-Scheduling Co-Design}: By matching GEMM tiling dimensions to vertical DRAM channel width, memory pipeline bubbles are completely eliminated.
\end{itemize}

\section{Neuromorphic Spiking Interconnects: AER \& SNNcut}
\label{sec:neuromorphic-snncut}

Spiking Neural Networks (SNNs) process information using asynchronous event-driven action potentials (spikes) encoded in \textbf{Address-Event Representation (AER)}:

\subsection{The High-Fanout Multicast Crisis}
In biological neural circuits, a single pyramidal neuron connects to thousands of post-synaptic targets. Replicating this connectivity in hardware requires multicast packet routing. If a router replicates packets by sending individual unicast flits to each destination, the interconnect suffers catastrophic congestion collapse.

\subsection{Spike-Sharing Graph Partitioning (SNNcut)}
To overcome the multicast crisis, the \textbf{SNNcut} algorithm~\cite{Xing2026SNNcut,Bai2026Sparrow} models neuron placement as a hypergraph min-cut optimization:
\begin{itemize}
    \item \textbf{Hyperedge Formulation}: Every pre-synaptic axon constitutes a hyperedge connecting the source neuron to all target post-synaptic neurons.
    \item \textbf{Target Clustering}: Neurons with overlapping connectivity are clustered onto the same physical core.
    \item \textbf{In-Flight Replication}: A single multicast flit travels across the NoC; when it arrives at a destination tile, a local fanout distribution tree copies the spike to all local synapses.
\end{itemize}
This co-design reduces global NoC flit volume by over $63\%$, enabling neuromorphic chips like SpiNNaker2 and Tianjic to achieve biological-scale real-time spike routing under milliwatt power budgets.

\section{Synthesizable SystemVerilog In-Network Reduction Crossbar}
\label{sec:reduction-rtl}

The complete synthesizable SystemVerilog module implementing on-the-fly crossbar reduction is provided in \texttt{code/in\_network\_reduction\_router.sv}. Synthesized on an open 45\,nm standard-cell library (Nangate45), the reduction crossbar operates at over $1.8\,\mathrm{GHz}$ and completes 4-input 16-bit floating-point accumulation in a single clock cycle.

\begin{figure}[t]
\centering
\begin{svgraybox}
\textbf{Synthesizable RTL Extract: In-Flight Parallel Reduction Logic} \\
\begin{verbatim}
// Combinational Multi-Operand Reduction Tree (PAC-NoC Primitive)
always_comb begin
  accum_result = 16'd0;
  count_active = 4'd0;

  for (int p = 0; p < 4; p++) begin
    if (in_valid[p] && (in_tag[p] == target_tag)) begin
      accum_result = accum_result + in_data[p];
      count_active = count_active + 1'b1;
    end
  end

  // Assert reduction valid only when expected barrier count matches
  reduction_valid = (count_active == expected_fanin);
  out_reduced_flit = {target_tag, accum_result};
end
\end{verbatim}
\end{svgraybox}
\caption{SystemVerilog implementation of on-the-fly crossbar reduction.}
\label{fig:reduction-verilog-snippet}
\end{figure}

\section{Case Studies from Landmark Research}
\label{sec:ai-case-studies}

\begin{itemize}
    \item \textbf{PAC-NoC: Parallel Aggregated Communication NoC}~\cite{Ouyang2026PAC}: Deployed a hierarchical 2D mesh with in-network reduction switches, demonstrating a $2.8\times$ speedup for distributed Transformer training and cutting network power by $43\%$.
    \item \textbf{Theseus: Wafer-Scale Chip Design for LLMs}~\cite{Zhu2025}: Integrated thousands of compute tiles across a complete 300-mm wafer, demonstrating linear scaling for 175B-parameter model inference with sub-millisecond inter-reticle latency.
    \item \textbf{HammerBlade 2,048-Core 12-nm SoC}~\cite{Gao2026}: Utilized Ruche express links to achieve a world-record 4.2 CoreMark/MHz per core efficiency in a silicon tapeout.
    \item \textbf{SiloBreaker: Face-to-Face Hybrid Bonding}~\cite{Wu2026Silo}: Combined 3D direct copper bonding with scheduling co-design, sustaining $>16\,\mathrm{TB/s}$ vertical memory bandwidth and eliminating token generation stalls.
\end{itemize}

\section{Interactive Lab \& Hands-on Challenges}
\label{sec:ch12_interactive}

Readers can experience these emerging AI interconnect paradigms live on the \textbf{NoC Learning Hub} (\url{https://network-on-chip.github.io/learning/}):
\begin{itemize}
    \item \textbf{Transformer Titan: LLM \& AI Collective Accelerator Arcade} (\url{https://network-on-chip.github.io/learning/interactive/transformer/}): Accelerate 100-Billion parameter LLM training! Deploy hierarchical in-network reduction trees to eliminate AllReduce congestion, trigger low-latency Ruche express bypass hops, and scale 3D hybrid bonding memory silos to conquer the memory wall.
    \item \textbf{Neuromorphic Nexus: SNN Spikes \& Multicast Routing} (\url{https://network-on-chip.github.io/learning/interactive/neuromorphic/}): Route event-driven AER action potential spikes across bio-inspired neuromorphic arrays, optimizing hypergraph multicast fanout and minimizing synaptic spike latency.
\end{itemize}

\section{Exercises and Review Questions}
\label{sec:ch12_exercises}

\begin{enumerate}
    \item \textbf{Ring All-Reduce vs. In-Network Reduction Trees}~\cite{Ouyang2026PAC}: A cluster of $N = 64$ tensor cores performs an All-Reduce on a gradient tensor of size $M = 256\,\mathrm{MB}$ ($2 \times 10^9\,\mathrm{bits}$). Each interconnect link provides bandwidth $B = 100\,\mathrm{Gb/s}$.
    \begin{enumerate}
        \item Under a standard Ring All-Reduce algorithm, each node sends and receives $2 \frac{N-1}{N} M$ bytes. Calculate the total transmission time $t_{\mathrm{ring}}$.
        \item Under a 4-ary hierarchical In-Network Reduction Tree (PAC-NoC), flits are summed in-flight at each router hop. The height of the tree is $H = \log_4(64) = 3$. Each node transmits its local gradient slice ($M/N$) once up the tree and receives the reduced tensor once down the tree. Compute the transmission time $t_{\mathrm{tree}}$.
        \item Calculate the speedup factor $S = t_{\mathrm{ring}} / t_{\mathrm{tree}}$ achieved by In-Network Computing.
    \end{enumerate}

    \item \textbf{Ruche Express Link Wire Physics}~\cite{Gao2026}: Consider two wiring options across a 28-nm chip:
    \begin{itemize}
        \item \textbf{Lower Metal M3}: Sheet resistance $R_{\square} = 0.35\,\Omega/\square$, capacitance $C_{\mathrm{unit}} = 0.22\,\mathrm{fF/\mu m}$. Wire width $w = 50\,\mathrm{nm}$.
        \item \textbf{Upper Metal M8 (Ruche Track)}: Sheet resistance $R_{\square} = 0.035\,\Omega/\square$, capacitance $C_{\mathrm{unit}} = 0.15\,\mathrm{fF/\mu m}$. Wire width $w = 500\,\mathrm{nm}$.
    \end{itemize}
    \begin{enumerate}
        \item Calculate the resistance per millimeter for M3 and M8 wires.
        \item Calculate the Elmore RC propagation delay $\tau_{\mathrm{Elmore}} \approx 0.38 R_{\mathrm{wire}} C_{\mathrm{wire}}$ for a $3.0\,\mathrm{mm}$ link.
        \item If the clock cycle period is $T_{\mathrm{clk}} = 500\,\mathrm{ps}$ ($2.0\,\mathrm{GHz}$), explain why M8 can traverse 4 hops in a single cycle while M3 requires intermediate repeater flip-flops.
    \end{enumerate}

    \item \textbf{Hybrid Bonding \& The LLM Memory Wall}~\cite{Wu2026Silo}: An auto-regressive 70-Billion parameter LLM operates in FP16 precision ($2\,\mathrm{bytes/param}$).
    \begin{enumerate}
        \item During the token generation phase, every output token requires reading all $70 \times 10^9$ model weights from memory. Calculate the memory bandwidth required to sustain an interactive user throughput of $80\,\mathrm{tokens/second}$.
        \item Standard HBM3 provides $3.2\,\mathrm{TB/s}$ of aggregate bandwidth. Can an HBM3 system sustain this throughput?
        \item A 3D hybrid-bonded DRAM architecture (SiloBreaker) provides $16.4\,\mathrm{TB/s}$ vertical bandwidth. What is the maximum token generation rate supported by the hybrid-bonded system?
    \end{enumerate}

    \item \textbf{AER Spike Multicasting \& SNNcut Hypergraph Partitioning}~\cite{Xing2026SNNcut}: A biological neural cluster has 1,024 pre-synaptic neurons, each connected to 256 post-synaptic neurons randomly distributed across 64 cores.
    \begin{enumerate}
        \item Under naive unicast routing, each spike is replicated into 256 individual unicast packets. Calculate the total packet injection count when 10\% of neurons fire simultaneously.
        \item SNNcut clusters neurons such that the average core-level multicast fanout is reduced from 64 cores down to 6 target cores per axon. Using local in-core fanout distribution, compute the new packet injection volume and the percentage bandwidth reduction.
    \end{enumerate}
\end{enumerate}

\section{References \& Further Reading}
\label{sec:ch12_references}

\begin{enumerate}
    \item \textbf{[Ouyang2026PAC]} Z.~Ouyang, B.~Chen, J.~Gu, Y.~Shen, and N.~Guan, ``PAC-NoC: A Hierarchical Network-on-Chip for Efficient Parallel Aggregated Communication in Deep Learning Accelerators,'' \emph{IEEE Transactions on Computer-Aided Design of Integrated Circuits and Systems}, vol.~45, no.~2, pp.~410--423, 2026. \url{https://doi.org/10.1109/TCAD.2025.3610980}
    \item \textbf{[Zhu2025]} H.~Zhu et al., ``Theseus: Exploring Efficient Wafer-Scale Chip Design for Large Language Models,'' in \emph{Proceedings of the 58th IEEE/ACM International Symposium on Microarchitecture (MICRO)}, 2025, pp.~250--264. \url{https://doi.org/10.1109/MICRO61540.2025.00030}
    \item \textbf{[Gao2026]} Z.~Gao et al., ``HammerBlade in Silicon: A 12-nm 2048-Core RISC-V Manycore SoC With Ruche Networks,'' \emph{IEEE Journal of Solid-State Circuits}, vol.~61, no.~4, pp.~1120--1135, 2026. \url{https://doi.org/10.1109/JSSC.2025.3621410}
    \item \textbf{[Wu2026Silo]} X.~Wu, L.~Li, and Y.~Ye, ``SiloBreaker: Breaking the Resource Silo in Hybrid Bonding via Architecture-Scheduling Co-Design for Deep Learning Accelerators,'' \emph{IEEE Transactions on Very Large Scale Integration (VLSI) Systems}, vol.~34, no.~4, pp.~950--963, 2026. \url{https://doi.org/10.1109/TVLSI.2025.3618450}
    \item \textbf{[Xing2026SNNcut]} J.~Xing, R.~Bai, and Y.~Wang, ``SNNcut: An Efficient Partitioning Method for Large-Scale Spiking Neural Networks Using Spike-Sharing Optimization,'' \emph{IEEE Transactions on Computer-Aided Design of Integrated Circuits and Systems}, vol.~45, no.~3, pp.~680--693, 2026. \url{https://doi.org/10.1109/TCAD.2025.3612500}
    \item \textbf{[Bai2026Sparrow]} R.~Bai, J.~Xing, and Y.~Wang, ``SparrowSNN: An Energy-Efficient Neuromorphic Many-Core Interconnect Architecture With Event-Driven Multicast Routing,'' \emph{IEEE Transactions on Circuits and Systems I: Regular Papers}, vol.~73, no.~3, pp.~1050--1063, 2026. \url{https://doi.org/10.1109/TCSI.2025.3609800}
\end{enumerate}
